% ============================================================
%  DATABASE DIAGRAMS & NORMALIZATION ANALYSIS
%  AI Cookbook System — ELTE BSc Thesis
%  Generated from ORM models in backend/models.py
% ============================================================


% ────────────────────────────────────────────────────────────
%  DIAGRAM 1 — Entity-Relationship Diagram
% ────────────────────────────────────────────────────────────

\begin{figure}[htbp]
    \centering
    \begin{plantuml}
    @startuml
    ' Entity-Relationship Diagram — AI Cookbook System

    skinparam classAttributeIconSize 0
    skinparam backgroundColor #FFFFFF
    skinparam classBackgroundColor #FFFFFF
    skinparam classHeaderBackgroundColor #FFFFFF
    skinparam packageBackgroundColor #FFFFFF
    skinparam classBorderColor #000000
    skinparam classArrowColor #000000
    skinparam packageBorderColor #000000
    skinparam fontColor #000000
    skinparam classFontColor #000000
    skinparam classAttributeFontColor #000000
    skinparam packageFontColor #000000
    skinparam classFontSize 13
    skinparam classAttributeFontSize 12
    skinparam packageFontSize 13
    skinparam classFontStyle bold
    skinparam nodesep 80
    skinparam ranksep 100
    skinparam padding 10

    hide circle
    top to bottom direction

    entity User {
        * id : INTEGER <<PK>>
        --
        * email : VARCHAR <<UNIQUE, NOT NULL>>
        * password : VARCHAR <<NOT NULL>>
        * preferences : TEXT <<NOT NULL, DEFAULT '[]'>>
        * allergies : TEXT <<NOT NULL, DEFAULT '[]'>>
        created_at : DATETIME(tz) <<DEFAULT now()>>
    }

    entity GeneratedRecipe {
        * id : INTEGER <<PK>>
        --
        * fingerprint : VARCHAR <<NOT NULL>>
        * title : VARCHAR <<NOT NULL>>
        * ingredients : TEXT <<NOT NULL>>
        * additional_ingredients : TEXT <<NOT NULL, DEFAULT '[]'>>
        * preferences : TEXT <<NOT NULL, DEFAULT '[]'>>
        * allergies : TEXT <<NOT NULL, DEFAULT '[]'>>
        * steps : TEXT <<NOT NULL>>
        * nutrition : TEXT <<NOT NULL>>
        * warnings : TEXT <<NOT NULL, DEFAULT '[]'>>
        image_url : VARCHAR
        image_path : VARCHAR
        image_cache_key : VARCHAR
        image_prompt : TEXT
        * is_saved : BOOLEAN <<NOT NULL, DEFAULT false>>
        saved_at : DATETIME(tz)
        created_at : DATETIME(tz) <<DEFAULT now()>>
        user_id : INTEGER <<FK>>
    }

    entity OcrExtraction {
        * id : INTEGER <<PK>>
        --
        * source_filename : VARCHAR <<NOT NULL>>
        * raw_text : TEXT <<NOT NULL>>
        * structured_nutrition : TEXT <<NOT NULL>>
        image_path : VARCHAR
        image_url : VARCHAR
        created_at : DATETIME(tz) <<DEFAULT now()>>
        * user_id : INTEGER <<FK, NOT NULL>>
    }

    User ||--o{ GeneratedRecipe : "generates"
    User ||--o{ OcrExtraction : "performs"

    @enduml
    \end{plantuml}
    \caption{Entity-relationship diagram of the AI Cookbook system database,
    showing all entities, attributes, primary keys, foreign keys,
    and cardinality constraints.}
    \label{fig:erd}
\end{figure}


% ────────────────────────────────────────────────────────────
%  DIAGRAM 2 — Relational Schema
% ────────────────────────────────────────────────────────────

\begin{figure}[htbp]
    \centering
    \begin{plantuml}
    @startuml
    ' Relational Schema Diagram — AI Cookbook System

    skinparam classAttributeIconSize 0
    skinparam backgroundColor #FFFFFF
    skinparam classBackgroundColor #FFFFFF
    skinparam classHeaderBackgroundColor #FFFFFF
    skinparam packageBackgroundColor #FFFFFF
    skinparam classBorderColor #000000
    skinparam classArrowColor #000000
    skinparam packageBorderColor #000000
    skinparam fontColor #000000
    skinparam classFontColor #000000
    skinparam classAttributeFontColor #000000
    skinparam packageFontColor #000000
    skinparam classFontSize 13
    skinparam classAttributeFontSize 12
    skinparam packageFontSize 13
    skinparam classFontStyle bold
    skinparam nodesep 80
    skinparam ranksep 100
    skinparam padding 10
    skinparam noteBackgroundColor #FFFFFF
    skinparam noteBorderColor #000000

    hide circle
    top to bottom direction

    class users <<Table>> {
        <<PK>> INTEGER id
        ..
        VARCHAR email <<UNIQUE, NOT NULL>>
        VARCHAR password <<NOT NULL>>
        TEXT preferences <<NOT NULL, DEFAULT '[]'>>
        TEXT allergies <<NOT NULL, DEFAULT '[]'>>
        DATETIME(tz) created_at <<DEFAULT now()>>
        ..
        **Indexes:**
        ix_users_id (id)
        ix_users_email (email, UNIQUE)
    }

    class generated_recipes <<Table>> {
        <<PK>> INTEGER id
        ..
        VARCHAR fingerprint <<NOT NULL>>
        VARCHAR title <<NOT NULL>>
        TEXT ingredients <<NOT NULL>>
        TEXT additional_ingredients <<NOT NULL, DEFAULT '[]'>>
        TEXT preferences <<NOT NULL, DEFAULT '[]'>>
        TEXT allergies <<NOT NULL, DEFAULT '[]'>>
        TEXT steps <<NOT NULL>>
        TEXT nutrition <<NOT NULL>>
        TEXT warnings <<NOT NULL, DEFAULT '[]'>>
        VARCHAR image_url
        VARCHAR image_path
        VARCHAR image_cache_key
        TEXT image_prompt
        BOOLEAN is_saved <<NOT NULL, DEFAULT false>>
        DATETIME(tz) saved_at
        DATETIME(tz) created_at <<DEFAULT now()>>
        <<FK>> INTEGER user_id -> users.id
        ..
        **Indexes:**
        ix_generated_recipes_id (id)
        ix_generated_recipes_fingerprint (fingerprint)
        ix_generated_recipes_image_cache_key (image_cache_key)
        ix_generated_recipes_user_saved_at (user_id, is_saved, saved_at)
        ix_generated_recipes_user_created_at (user_id, created_at)
    }

    class ocr_extractions <<Table>> {
        <<PK>> INTEGER id
        ..
        VARCHAR source_filename <<NOT NULL>>
        TEXT raw_text <<NOT NULL>>
        TEXT structured_nutrition <<NOT NULL>>
        VARCHAR image_path
        VARCHAR image_url
        DATETIME(tz) created_at <<DEFAULT now()>>
        <<FK>> INTEGER user_id <<NOT NULL>> -> users.id
        ..
        **Indexes:**
        ix_ocr_extractions_id (id)
        ix_ocr_extractions_user_created_at (user_id, created_at)
    }

    users "1" --o "0..N" generated_recipes : "user_id"
    users "1" --o "1..N" ocr_extractions : "user_id"

    note "All tables satisfy BCNF:\n- No partial dependencies (all PKs are single-column surrogate keys)\n- No transitive dependencies among non-key attributes\n- All determinants are candidate keys\n\nIntentional denormalization:\n- preferences, allergies, ingredients, warnings, steps,\n  nutrition, and structured_nutrition store JSON-serialized\n  TEXT to avoid excessive join complexity for\n  variable-length lists returned as atomic API payloads." as NormNote

    @enduml
    \end{plantuml}
    \caption{Relational schema derived from the ER model,
    showing all tables with primary keys, foreign keys,
    integrity constraints, and inter-table relationships.}
    \label{fig:relational-schema}
\end{figure}


% ────────────────────────────────────────────────────────────
%  NORMALIZATION ANALYSIS
% ────────────────────────────────────────────────────────────

The database schema of the AI Cookbook system consists of three
relations: \texttt{users}, \texttt{generated\_recipes}, and
\texttt{ocr\_extractions}. Each relation uses a single-column
surrogate integer primary key (\texttt{id}), which is
auto-incremented by the database engine. The following analysis
examines the normalization level achieved by this schema and
identifies any intentional deviations from strict normal forms.

Regarding First Normal Form, every column in all three relations
stores a single atomic value at the SQL type level. There are no
repeating groups in the relational sense: each row in
\texttt{generated\_recipes} represents exactly one recipe, and each
row in \texttt{ocr\_extractions} represents exactly one extraction
event. The \texttt{email} column in the \texttt{users} table stores
a single string value constrained to be unique across the relation,
and the \texttt{password} column stores a single hashed string.
Columns such as \texttt{preferences} and \texttt{allergies} in the
\texttt{users} table, as well as \texttt{ingredients},
\texttt{additional\_ingredients}, \texttt{steps}, \texttt{nutrition},
\texttt{warnings}, and \texttt{structured\_nutrition} in their
respective tables, store JSON-serialized text. While these columns
contain structured data within a single TEXT field, the database
engine treats each value as an indivisible string, and therefore the
schema satisfies 1NF under the standard relational definition.

Regarding Second Normal Form, every non-key attribute in each
relation is fully functionally dependent on the entire primary key.
Since all three primary keys consist of a single surrogate column
(\texttt{id}), partial dependencies are structurally impossible.
For instance, in the \texttt{generated\_recipes} table, every
attribute---\texttt{fingerprint}, \texttt{title},
\texttt{ingredients}, \texttt{is\_saved}, \texttt{user\_id}, and
all others---depends on the single primary key \texttt{id} and
cannot depend on a proper subset of it. The same reasoning applies
to \texttt{users.id} and \texttt{ocr\_extractions.id}. Therefore,
all three relations satisfy 2NF.

Regarding Third Normal Form and Boyce--Codd Normal Form, there are
no transitive dependencies among non-key attributes. In the
\texttt{users} relation, \texttt{email} is a candidate key (enforced
by the unique constraint), and the remaining attributes
\texttt{password}, \texttt{preferences}, \texttt{allergies}, and
\texttt{created\_at} each depend directly on \texttt{id} (or
equivalently on \texttt{email}) without any intermediate non-key
determinant. In the \texttt{generated\_recipes} relation,
\texttt{user\_id} is a foreign key referencing \texttt{users.id},
but it is not a determinant of any other non-key attribute within
the same relation; it merely records ownership. The
\texttt{fingerprint} column serves as a content hash for
deduplication lookups but does not functionally determine any other
column, since the same fingerprint could theoretically appear in
different user contexts. All remaining columns---recipe content
fields, image metadata, and timestamps---depend solely on the
primary key. The \texttt{ocr\_extractions} relation follows the same
pattern: \texttt{source\_filename}, \texttt{raw\_text},
\texttt{structured\_nutrition}, image fields, and
\texttt{created\_at} all depend exclusively on \texttt{id}. Since
every determinant in all three relations is a superkey, the schema
satisfies BCNF.

The schema does contain intentional denormalization in several TEXT
columns that store JSON-serialized arrays or objects rather than
being decomposed into separate child tables. Specifically, the
\texttt{preferences} and \texttt{allergies} columns in both
\texttt{users} and \texttt{generated\_recipes} store serialized JSON
arrays of strings (for example, \texttt{["vegetarian",
"low-carb"]}). The \texttt{ingredients},
\texttt{additional\_ingredients}, \texttt{steps}, \texttt{warnings},
and \texttt{nutrition} columns in \texttt{generated\_recipes}, as
well as \texttt{structured\_nutrition} in
\texttt{ocr\_extractions}, similarly store JSON-encoded structured
data. A strictly normalized design would decompose each of these
into a separate relation---for example, a
\texttt{recipe\_ingredients} junction table with one row per
ingredient per recipe. However, this denormalization is a deliberate
design choice driven by the application's access patterns: every
API endpoint that retrieves a recipe or extraction returns these
fields as complete JSON arrays in a single response payload, and no
query in the application ever filters, joins, or aggregates on
individual elements within these fields. Decomposing them would
introduce significant join overhead for every read operation while
providing no query-time benefit, since the application never needs
to answer questions such as ``which recipes contain ingredient X?''
at the database level. The tradeoff is that these fields cannot be
efficiently queried or indexed at the element level by the database
engine, but this limitation is acceptable because such queries are
handled by the separate Retrieval-Augmented Generation (RAG)
subsystem, which uses a ChromaDB vector store rather than relational
SQL queries for semantic ingredient and recipe search.
